\section{Coefficient identity, copy weights, and the block-diagonal gauge}
\label{sec:sector_bnt_full_basis_global_gauge}

The bijective matching theorem of Chapter~\ref{ch:bnt}
(Theorem~\ref{thm:sector_bnt_bijective_match_of_sameMPV}) gives the sector
phases $\zeta_k$ and block gauges $X_k$. The coefficient comparison of
\cite[Appendix MPV proof, lines~1187--1188]{Cirac2016MPDO_arXiv} then turns the
equality of total MPV families into the per-sector coefficient identity
$c_N^{(\beta(k))}(P)=\zeta_k^N\,c_N^{(k)}(Q)$. Neither
step uses a per-sector unit-modulus hypothesis.

\begin{lemma}[Unit phase from matched normalized BNT blocks]%
    \label{lem:sector_bnt_norm_phase_of_matched_mpv}
    \lean{MPSTensor.IsBNTCanonicalForm.norm_phase_of_matched_mpv}
    \leanok
    \uses{def:sector_bnt_cf}
    Let $P$ and $Q$ be BNT canonical forms. If, for every positive length $N$
    and word $\sigma$, a sector $P_j$ and a sector $Q_k$ have MPV families
    related by
    \begin{align}
        V^{(N)}(Q_k)_\sigma
        &= \zeta^N V^{(N)}(P_j)_\sigma,
        \label{eq:ft_matched_mpv_phase_norm}
    \end{align}
    then $|\zeta|=1$.
\end{lemma}

\begin{proof}\leanok
    The BNT normalization gives self-overlap limits of modulus $1$ for both
    sectors. The MPV identity (\ref{eq:ft_matched_mpv_phase_norm}) gives, for
    every positive $N$,
    \begin{align}
        O_{Q_kQ_k}(N)
        &= |\zeta|^{2N}\,O_{P_jP_j}(N).
        \notag
    \end{align}
    Since $O_{Q_kQ_k}(N)\to1$ and $O_{P_jP_j}(N)\to1$, taking limits forces
    $|\zeta|^{2N}\to1$ and hence $|\zeta|=1$.
\end{proof}

\begin{theorem}[Proportional sector matching with unit phases]%
    \label{thm:sector_bnt_proportional_sector_match_witnesses}
    \lean{MPSTensor.ft_sector_bnt_proportional_sector_match_witnesses}
    \leanok
    \uses{def:sector_bnt_cf,
        def:gauge_phase_equiv,
        thm:gauge_phase_mpv_scaling,
        lem:sector_bnt_norm_phase_of_matched_mpv}
    Let $P$ and $Q$ be BNT canonical forms whose total tensors generate
    eventually nonzero proportional MPV families.
    Then there are a bijection $\beta:J(Q)\simeq J(P)$, bond-dimension
    equalities $\dim P_{\beta(k)}=\dim Q_k$, scalars
    $\zeta_k\in\C$, and block gauges $X_k$ such that $|\zeta_k|=1$ and, for
    every sector $k$ and physical index $i$,
    \begin{align}
        Q_k^i
        &= \zeta_k X_k P_{\beta(k)}^i X_k^{-1}.
        \label{eq:ft_gauge_relation}
    \end{align}
    Diagrammatically, the gauge relation (\ref{eq:ft_gauge_relation}) reads
    \begin{center}
        % Q_k^i = zeta_k X_k P_{beta(k)}^i X_k^{-1}: the gauge X_k and its
        % inverse ride the virtual bond.  Both sides expose one virtual
        % index at each end and one physical leg, so the two sides carry
        % the same boundary signature
        % (virtual-west=1 | virtual-east=1 | physical-up=1).
        \begin{tenkz}[physical=up]
            \tn[up=$i$]{Q_k}
        \end{tenkz}
        $ = \zeta_k$
        \begin{tenkz}[physical=up]
            \tnX{X_k} & \tn[up=$i$]{P_{\beta(k)}} & \tnX{X_k^{-1}}
        \end{tenkz}
    \end{center}
    with every open virtual end a matrix index.
    Equivalently, for every positive length $N$ and word $\sigma$,
    \begin{align}
        V^{(N)}(Q_k)_\sigma
        &= \zeta_k^N V^{(N)}(P_{\beta(k)})_\sigma.
        \notag
    \end{align}
    This is \cite[Appendix MPV theorem statement, lines~1167--1170, and proof
        line~1182]{Cirac2016MPDO_arXiv}.
    % Per-sector unit-modulus restriction eliminated by the exact matcher;
    % closure recorded in
    % docs/paper-gaps/cpsv16_global_vs_persector_unit_witness.tex.
\end{theorem}

\begin{proof}\leanok\uses{lem:sector_bnt_norm_phase_of_matched_mpv}
    Apply the proportional sector-matching theorem to obtain the bijection and
    dimension-compatible gauge-phase equivalences. Choosing the gauges and
    scalars gives $Q_k^i=\zeta_kX_kP_{\beta(k)}^iX_k^{-1}$, and therefore
    \begin{align}
        O_{Q_kQ_k}(N)
        &= |\zeta_k|^{2N}O_{P_{\beta(k)}P_{\beta(k)}}(N).
        \notag
    \end{align}
    The normalized self-overlap limits give
    $O_{Q_kQ_k}(N)\to1$ and $O_{P_{\beta(k)}P_{\beta(k)}}(N)\to1$, so
    $|\zeta_k|^{2N}\to1$ and hence $|\zeta_k|=1$. The MPV scalar-power
    identity follows from the same gauge equation (\ref{eq:ft_gauge_relation}):
    $V^{(N)}(Q_k)_\sigma=\zeta_k^NV^{(N)}(P_{\beta(k)})_\sigma$.
\end{proof}

\begin{lemma}[Full-basis coefficient identity from matched sectors]%
    \label{lem:sector_bnt_coeff_identity_via_global_gauge}
    \lean{MPSTensor.coeff_identity_via_global_gauge}
    \leanok
    \uses{def:sector_bnt_cf,
        def:gauge_phase_equiv}
    Let $P$ and $Q$ be BNT canonical forms with the same MPV family at every
    positive length, matched by a bijection $\beta:J(Q)\simeq J(P)$ with
    bond-dimension equalities and a
    gauge-phase equivalence between $Q_k$ and $P_{\beta(k)}$ for every $k$. Then
    each sector carries a unit-modulus phase $\zeta_k$, and eventually
    \begin{align}
        c_N^{(\beta(k))}(P)
        &= \zeta_k^N\,c_N^{(k)}(Q).
        \label{eq:ft_coeff_identity}
    \end{align}
    This is the equal-MPV coefficient comparison of
    \cite[Appendix MPV proof, lines~1187--1188]{Cirac2016MPDO_arXiv}.
\end{lemma}

\begin{proof}\leanok
    \uses{lem:sector_bnt_norm_phase_of_matched_mpv,
        lem:sector_bnt_coeff_identity_via_matched_mpv_phase}
    For each matched sector, the gauge-phase equivalence gives
    $V^{(N)}(Q_k)_\sigma=\zeta_k^NV^{(N)}(P_{\beta(k)})_\sigma$.
    Lemma~\ref{lem:sector_bnt_norm_phase_of_matched_mpv} gives
    $|\zeta_k|=1$. Applying
    Lemma~\ref{lem:sector_bnt_coeff_identity_via_matched_mpv_phase} to the
    same phases gives
    $c_N^{(\beta(k))}(P)=\zeta_k^Nc_N^{(k)}(Q)$ for all sufficiently
    large $N$.
\end{proof}

\begin{lemma}[Full-basis coefficient identity from matched phases]%
    \label{lem:sector_bnt_coeff_identity_via_matched_mpv_phase}
    \lean{MPSTensor.coeff_identity_via_matched_mpv_phase}
    \leanok
    \uses{def:sector_bnt_cf}
    Let $P$ be a BNT canonical form, let $Q$ be a sector decomposition, and
    assume $P_{\mathrm{tot}}$ and $Q_{\mathrm{tot}}$ have the same MPV family
    at every positive length.
    Suppose that there is a bijection $\beta:J(Q)\simeq J(P)$ and scalars
    $\zeta_k \in \C$ such that, for every $k \in J(Q)$, every positive length
    $N$, and every spin configuration $\sigma$ of length $N$,
    \begin{align}
        V^{(N)}(Q_k)_\sigma
        &= \zeta_k^N\,V^{(N)}(P_{\beta(k)})_\sigma.
        \label{eq:ft_phase_relation_hyp}
    \end{align}
    Then for every $k \in J(Q)$ there is $N_0$ such that, for all $N>N_0$,
    $c_N^{(\beta(k))}(P)=\zeta_k^N\,c_N^{(k)}(Q)$.
    This is the coefficient-comparison step of
    \cite[Appendix MPV proof, lines~1187--1188]{Cirac2016MPDO_arXiv} in full-basis
    form.
\end{lemma}

\begin{proof}\leanok
    \uses{lem:eventual_coeff_eq_of_eventual_li}
    For every $N$, expand the equality of total MPVs as
    \begin{align}
        \sum_{j\in J(P)} c_N^{(j)}(P)\ket{V^{(N)}(P_j)}
        &=
        \sum_{k\in J(Q)} c_N^{(k)}(Q)\ket{V^{(N)}(Q_k)}.
        \notag
    \end{align}
    The phase relation (\ref{eq:ft_phase_relation_hyp}) gives
    $\ket{V^{(N)}(Q_k)}=\zeta_k^N\ket{V^{(N)}(P_{\beta(k)})}$, so reindexing
    the $Q$-sum by $\beta$ onto the $P$-basis yields
    \begin{align}
        \sum_{j\in J(P)} c_N^{(j)}(P)\ket{V^{(N)}(P_j)}
        &=
        \sum_{j\in J(P)}
        \zeta_{\beta^{-1}(j)}^N
        c_N^{(\beta^{-1}(j))}(Q)\ket{V^{(N)}(P_j)}.
        \notag
    \end{align}
    For all sufficiently large $N$ the $P$-basis MPV family
    $\{\ket{V^{(N)}(P_j)}\}_{j\in J(P)}$ is linearly independent, by the BNT
    eventual linear independence of $P$. Lemma~\ref{lem:eventual_coeff_eq_of_eventual_li}
    then extracts the coefficient equalities
    $c_N^{(\beta(k))}(P)=\zeta_k^Nc_N^{(k)}(Q)$ from this eventual linear
    independence, for all sufficiently large $N$.
\end{proof}

The coefficient identity (\ref{eq:ft_coeff_identity}) is now converted, sector
by sector, into a matching of the individual copy weights.

\begin{definition}[Copy-weight matching for matched BNT sectors]%
    \label{def:sector_bnt_copy_weight_matching}
    \lean{MPSTensor.SectorBNTCopyWeightMatching}
    \leanok
    Let the sectors of two decompositions $P$ and $Q$ be matched by
    $\beta:J(Q)\simeq J(P)$, with sector phases $\zeta_k$. A copy-weight matching
    consists of copy permutations $\tau_k$ such that, for all sectors $k$ and
    copies $q$, $\nu_{k,q}=\zeta_k^{-1}\mu_{\beta(k),\tau_k(q)}$.
    This is the copy-level conclusion of
    \cite[Appendix MPV proof, line~1188]{Cirac2016MPDO_arXiv}.
\end{definition}

\begin{lemma}[Copy-weight matching from matched-sector coefficient identities]%
    \label{lem:sector_bnt_copy_weight_matching_of_coeff_identity}
    \lean{MPSTensor.SectorBNTCopyWeightMatching.of_coeff_identity}
    \leanok
    \uses{def:sector_bnt_copy_weight_matching,
        lem:sector_bnt_matched_sector_weight_equiv}
    Let the sectors of $P$ and $Q$ be matched by $\beta:J(Q)\simeq J(P)$, with
    nonzero sector phases $\zeta_k$. If the matched-sector coefficient identities
    \begin{align}
        c_N^{(\beta(k))}(P)
        &= \zeta_k^N\,c_N^{(k)}(Q)
        \label{eq:ft_sector_coeff_identity_hyp}
    \end{align}
    hold eventually, then the phases $\zeta_k$ determine a copy-weight matching
    $\nu_{k,q}=\zeta_k^{-1}\mu_{\beta(k),\tau_k(q)}$. This is, sector by sector,
    \cite[Appendix MPV proof, line~1188]{Cirac2016MPDO_arXiv}.
\end{lemma}

\begin{proof}\leanok
    Eventually the coefficient identity
    (\ref{eq:ft_sector_coeff_identity_hyp}) in sector $k$ reads
    \begin{align}
        \sum_r \mu_{\beta(k),r}^{\,N}
        &= \zeta_k^N\sum_q\nu_{k,q}^{\,N}
        = \sum_q(\zeta_k\nu_{k,q})^{N}.
        \notag
    \end{align}
    Newton--Girard finite power-sum comparison
    (Lemma~\ref{lem:sector_bnt_matched_sector_weight_equiv}) applies sector by
    sector, producing a copy permutation $\tau_k$ with
    $\mu_{\beta(k),\tau_k(q)}=\zeta_k\nu_{k,q}$; since $\zeta_k\ne0$ this is
    $\nu_{k,q}=\zeta_k^{-1}\mu_{\beta(k),\tau_k(q)}$.
\end{proof}

The matched block gauges and copy weights now combine into a single
block-diagonal gauge.

\begin{definition}[Flattened block-diagonal gauge]%
    \label{def:global_gauge_of_blocks}
    \lean{MPSTensor.blockDiagonalGL, MPSTensor.globalGaugeOfBlocks}
    \leanok
    \uses{def:block_diagonal_tensor}
    Given per-block gauges $X_k \in \GL_{D_k}(\C)$, form the block-diagonal
    element $\bigoplus_k X_k$ and then reindex the dependent block-coordinate
    space with the standard bond index
    $\{0,\ldots,\sum_k D_k-1\}$ of the assembled tensor. The result is an
    element of $\GL_{\sum_k D_k}(\C)$.
\end{definition}

\begin{lemma}[Unitary block-diagonal gauges]
    \label{lem:unitary_global_gauge_of_blocks}
    \lean{MPSTensor.unitaryGL,
        MPSTensor.globalGaugeOfBlocks_unitaryGL_mem}
    \leanok
    \uses{def:global_gauge_of_blocks}
    A unitary matrix $U_k$ defines an invertible block gauge with inverse
    $U_k^\dagger$. If every block gauge is unitary, then the flattened
    block-diagonal gauge $U=\bigoplus_k U_k$ is unitary.
\end{lemma}
\begin{proof}\leanok
    The products are computed blockwise:
    \begin{align}
        U U^\dagger
        &= \bigoplus_k U_kU_k^\dagger
        = \bigoplus_k\mathbb 1
        = \mathbb 1.
        \notag
    \end{align}
\end{proof}

\begin{lemma}[Direct-sum conjugation by the flattened gauge]%
    \label{lem:tensor_from_blocks_globalGauge_conj}
    \lean{MPSTensor.toTensorFromBlocks_eq_globalGaugeOfBlocks_conj,
        MPSTensor.gaugeEquiv_toTensorFromBlocks_of_blockConj}
    \leanok
    \uses{def:global_gauge_of_blocks, def:block_diagonal_tensor, def:gauge_equiv}
    If, for every block $k$ and physical index $i$,
    $B_k^i=X_kA_k^iX_k^{-1}$, then the weighted direct sums satisfy
    \begin{align}
        \bigoplus_k \mu_k B_k^i
        &=
        X\left(\bigoplus_k \mu_k A_k^i\right)X^{-1},
        \notag
    \end{align}
    where $X$ is the flattened block-diagonal gauge from
    Definition~\ref{def:global_gauge_of_blocks}. Hence the two assembled
    tensors are gauge equivalent.
\end{lemma}

\begin{proof}\leanok
    Multiplication of block-diagonal matrices is computed blockwise. The scalar
    weight $\mu_k$ commutes with the conjugation on the $k$th block, and the
    reindexing from dependent block coordinates to the flattened bond index
    preserves products and inverses.
\end{proof}

\begin{theorem}[Global gauge from matched sectors and copy weights]%
    \label{thm:sector_bnt_global_gauge_of_matched_weights}
    \lean{MPSTensor.sector_bnt_global_gauge_of_matched_weights}
    \leanok
    \uses{lem:tensor_from_blocks_globalGauge_conj}
    Suppose that the sectors of two sector decompositions $P$ and $Q$ are matched
    by a bijection $\beta:J(Q)\simeq J(P)$, with bond-dimension equalities,
    copy permutations $\tau_k$, nonzero phases $\zeta_k$, and block gauges
    $X_k$. If, for every sector $k$, copy $q$, and physical index $i$,
    \begin{align}
        Q_k^i
        &= \zeta_k\,X_k P_{\beta(k)}^i X_k^{-1},
        \notag\\
        \nu_{k,q}
        &= \zeta_k^{-1} \mu_{\beta(k),\tau_k(q)},
        \label{eq:ft_matched_weights_hypotheses}
    \end{align}
    then, in the flattened $Q$-coordinates,
    \begin{align}
        Q_{\mathrm{flat}}^i
        &= X\,P_{\mathrm{matched}}^i\,X^{-1},
        \notag\\
        X
        &= \bigoplus_k(\mathbf{1}_{r_k}\otimes X_k).
        \notag
    \end{align}
    This is the direct-sum gauge construction of
    \cite[Appendix MPV proof, lines~1189--1192]{Cirac2016MPDO_arXiv}, separated from
    the coefficient-comparison step that supplies the weight identities.
\end{theorem}

\begin{proof}\leanok
    For each matched copy,
    \begin{align}
        \nu_{k,q}Q_k^i
        &= (\zeta_k^{-1}\mu_{\beta(k),\tau_k(q)})
        (\zeta_kX_kP_{\beta(k)}^iX_k^{-1})
        = \mu_{\beta(k),\tau_k(q)}X_kP_{\beta(k)}^iX_k^{-1}.
        \notag
    \end{align}
    Thus every flattened weighted block is conjugated by its sector gauge
    $X_k$. Lemma~\ref{lem:tensor_from_blocks_globalGauge_conj} combines these
    block conjugacies into the direct-sum gauge
    $X=\bigoplus_k(\mathbf{1}_{r_k}\otimes X_k)$.
\end{proof}

\begin{theorem}[Global gauge from a copy-weight matching]%
    \label{thm:sector_bnt_global_gauge_of_copy_weight_matching}
    \lean{MPSTensor.sector_bnt_global_gauge_of_copy_weight_matching}
    \leanok
    \uses{def:sector_bnt_copy_weight_matching,
        thm:sector_bnt_global_gauge_of_matched_weights}
    Suppose that the sectors of $P$ and $Q$ are matched, with
    bond-dimension equalities, nonzero phases, and block gauges satisfying
    $Q_k^i=\zeta_k X_kP_{\beta(k)}^iX_k^{-1}$.
    If the same phases also give a copy-weight matching, then the flattened
    $Q$-tensor is obtained from the matched-coordinate $P$-tensor by the
    direct-sum gauge $X=\bigoplus_k(\mathbf{1}_{r_k}\otimes X_k)$.
    This is only a reformulation of
    Theorem~\ref{thm:sector_bnt_global_gauge_of_matched_weights}.
\end{theorem}

\begin{proof}\leanok
    The copy-weight matching supplies
    $\nu_{k,q}=\zeta_k^{-1}\mu_{\beta(k),\tau_k(q)}$. Together with
    $Q_k^i=\zeta_kX_kP_{\beta(k)}^iX_k^{-1}$, these are exactly the two
    hypotheses (\ref{eq:ft_matched_weights_hypotheses}) of
    Theorem~\ref{thm:sector_bnt_global_gauge_of_matched_weights}.
\end{proof}

We can now state the proportional global gauge, first from a copy-weight
matching and then from the coefficient identities that produce one.

\begin{theorem}[Conditional proportional global gauge from a copy-weight matching]%
    \label{thm:sector_bnt_proportional_global_gauge_of_copy_weight_matching}
    \lean{MPSTensor.ft_sector_bnt_proportional_global_gauge_of_copy_weight_matching}
    \leanok
    \uses{def:sector_bnt_copy_weight_matching,
        thm:sector_bnt_proportional_sector_match_witnesses,
        thm:sector_bnt_global_gauge_of_copy_weight_matching}
    Let $P$ and $Q$ be BNT canonical forms whose total tensors generate
    eventually nonzero proportional MPV families.
    Then the proportional sector-matching theorem gives $\beta$, the
    bond-dimension equalities, the unit phases $\zeta_k$, and the block gauges
    $X_k$. If the remaining proportional coefficient comparison supplies a
    copy-weight matching for these same phases, then the flattened
    $Q$-tensor is obtained from the matched-coordinate $P$-tensor by the
    direct-sum gauge $X=\bigoplus_k(\mathbf{1}_{r_k}\otimes X_k)$.
    % Per-sector unit-modulus restriction eliminated by the exact matcher;
    % closure recorded in
    % docs/paper-gaps/cpsv16_global_vs_persector_unit_witness.tex.
\end{theorem}

\begin{proof}\leanok
    The proportional sector-matching theorem supplies $\beta$, $X_k$, and
    unit phases $\zeta_k$ satisfying
    $Q_k^i=\zeta_kX_kP_{\beta(k)}^iX_k^{-1}$. Since $|\zeta_k|=1$, each
    $\zeta_k$ is nonzero. A copy-weight matching gives
    $\nu_{k,q}=\zeta_k^{-1}\mu_{\beta(k),\tau_k(q)}$, so
    Theorem~\ref{thm:sector_bnt_global_gauge_of_copy_weight_matching} gives the
    direct-sum gauge.
\end{proof}

\begin{theorem}[Conditional proportional global gauge from coefficient identities]%
    \label{thm:sector_bnt_proportional_global_gauge_of_coeff_identity}
    \lean{MPSTensor.ft_sector_bnt_proportional_global_gauge_of_coeff_identity}
    \leanok
    \uses{lem:sector_bnt_copy_weight_matching_of_coeff_identity,
        thm:sector_bnt_proportional_global_gauge_of_copy_weight_matching}
    Let $P$ and $Q$ be BNT canonical forms whose total tensors generate
    eventually nonzero proportional MPV families.
    Then there are a bijection $\beta:J(Q)\simeq J(P)$, bond-dimension
    equalities, unit-modulus phases $\zeta_k$, and block gauges $X_k$ such that
    $Q_k^i=\zeta_k X_kP_{\beta(k)}^iX_k^{-1}$ for every sector $k$ and physical
    index $i$. If these same phases satisfy the matched-sector coefficient
    identities $c_N^{(\beta(k))}(P)=\zeta_k^N\,c_N^{(k)}(Q)$ for all
    sufficiently large $N$, then they determine a copy-weight matching and
    hence the flattened $Q$-tensor is obtained from the matched-coordinate
    $P$-tensor by the direct-sum gauge
    $X=\bigoplus_k(\mathbf{1}_{r_k}\otimes X_k)$.
    % Per-sector unit-modulus restriction eliminated by the exact matcher;
    % closure recorded in
    % docs/paper-gaps/cpsv16_global_vs_persector_unit_witness.tex.
\end{theorem}

\begin{proof}\leanok
    The sector-matching theorem gives $\beta$, $X_k$, and $|\zeta_k|=1$ with
    $Q_k^i=\zeta_kX_kP_{\beta(k)}^iX_k^{-1}$. Since $\zeta_k\ne0$, the
    coefficient identities
    $c_N^{(\beta(k))}(P)=\zeta_k^Nc_N^{(k)}(Q)$ produce
    $\nu_{k,q}=\zeta_k^{-1}\mu_{\beta(k),\tau_k(q)}$ by
    Lemma~\ref{lem:sector_bnt_copy_weight_matching_of_coeff_identity}.
    Theorem~\ref{thm:sector_bnt_proportional_global_gauge_of_copy_weight_matching}
    then applies to this copy-weight matching.
\end{proof}

The coefficient-identity hypothesis of
Theorem~\ref{thm:sector_bnt_proportional_global_gauge_of_coeff_identity} is the
scalar-free identity (\ref{eq:ft_coeff_identity}). From the proportional MPV
hypothesis alone one obtains the same identity carrying the global length-scalar
of the proportionality, recorded next.

\begin{lemma}[Proportional matched-phase coefficient identity, scalar-threaded]%
    \label{lem:sector_bnt_coeff_identity_via_matched_mpv_phase_proportional}
    \lean{MPSTensor.coeff_identity_via_matched_mpv_phase_proportional}
    \leanok
    \uses{def:sector_bnt_cf, lem:eventual_coeff_eq_of_eventual_li}
    Let $P$ be a BNT canonical form and $Q$ a sector decomposition whose total
    tensors generate eventually nonzero proportional MPV families. Suppose a
    bijection $\beta:J(Q)\simeq J(P)$ and scalars $\zeta_k$ satisfy the matched
    MPV phase identities (\ref{eq:ft_phase_relation_hyp}) for every sector $k$.
    Then there is a single eventually-nonzero scalar sequence $c_N$ and a
    threshold $N_0$ such that, for all $N>N_0$ and every $k\in J(Q)$,
    \begin{align}
        c_N^{(\beta(k))}(P)
        &= c_N\,\zeta_k^N\,c_N^{(k)}(Q).
        \label{eq:ft_coeff_identity_threaded}
    \end{align}
    The scalar $c_N$ is the same for every sector $k$: it is the single global
    proportionality scalar of \cite[Theorem~\texttt{thm1}, lines~1167--1170]%
    {Cirac2016MPDO_arXiv} at length $N$.
\end{lemma}

\begin{proof}\leanok
    For each $N$ on the proportionality tail, expand both total MPVs in their
    sector bases and use the global proportionality scalar $c_N$:
    \begin{align}
        \sum_{j\in J(P)} c_N^{(j)}(P)\ket{V^{(N)}(P_j)}
        &=
        c_N\sum_{k\in J(Q)} c_N^{(k)}(Q)\ket{V^{(N)}(Q_k)}.
        \notag
    \end{align}
    The matched phase identities (\ref{eq:ft_phase_relation_hyp}) give
    $\ket{V^{(N)}(Q_k)}=\zeta_k^N\ket{V^{(N)}(P_{\beta(k)})}$; because every
    sector is matched, reindexing the $Q$-sum by $\beta$ rewrites the right-hand
    side entirely in the $P$-basis. The BNT eventual linear independence of the
    $P$-basis and Lemma~\ref{lem:eventual_coeff_eq_of_eventual_li} then extract
    the coefficient identities (\ref{eq:ft_coeff_identity_threaded}) for all
    sufficiently large $N$.
\end{proof}

\begin{remark}[Residual length-scalar control]%
    \label{rem:proportional_length_scalar_gap}
    The identity (\ref{eq:ft_coeff_identity_threaded}) discharges the hypothesis
    of Theorem~\ref{thm:sector_bnt_proportional_global_gauge_of_coeff_identity}
    --- equivalently (\ref{eq:ft_coeff_identity}) --- precisely when
    $c_N\equiv1$ on the proportionality tail. More generally, if $c_N=\xi^N$ for one fixed phase
    $\xi$, then the coefficient identity uses the phases $\xi\zeta_k$, whereas
    the matched block relation still uses $\zeta_k$. Copy-weight recovery and
    block-gauge assembly therefore leave the common factor $\xi^{-1}$ and yield
    conjugacy up to that phase, not the exact conjugacy asserted by the current
    conditional theorem. Establishing such a pure phase-power form for the
    global proportionality scalar is a separate length-scalar control statement
    requiring peripheral-spectrum or almost-periodicity analysis of the bounded
    power sums $c_N^{(j)}(\cdot)$; it is not derivable from
    (\ref{eq:ft_coeff_identity_threaded}) alone. This distinction is consistent
    with the sectorwise gauge-phase conclusion of
    \cite[Theorem~\texttt{thm1}]{Cirac2016MPDO_arXiv}. Until an appropriate
    full-tensor theorem allowing the residual common phase is supplied, the
    unconditional full-tensor proportional gauge equivalence is not stated; the
    unconditional per-normal-tensor proportional fundamental theorem is
    Theorem~\ref{thm:cpgsv_multiblock_ft_source}.
    % Paper-gap note:
    % docs/paper-gaps/cpsv16_proportional_length_scalar_gap.tex.
\end{remark}

In the equal-MPV case all ingredients are unconditional.

\begin{lemma}[Equal-MPV matched copy-weight witnesses]%
    \label{lem:sector_bnt_equal_matched_copy_weight_witnesses}
    \lean{MPSTensor.ft_sector_bnt_equal_matched_copy_weight_witnesses}
    \leanok
    \uses{def:sector_bnt_cf}
    Let $P$ and $Q$ be BNT canonical forms whose total tensors have the same
    MPV family at every positive length. Then there are a bijection
    $\beta:J(Q)\simeq J(P)$, bond-dimension equalities, unit-modulus phases
    $\zeta_k$, and block gauges $X_k$ such that
    $Q_k^i=\zeta_k X_kP_{\beta(k)}^iX_k^{-1}$.
    For the same phases there is a copy-weight matching: after a permutation
    of the copies inside each matched sector,
    $\nu_{k,q}=\zeta_k^{-1}\mu_{\beta(k),\tau_k(q)}$.
\end{lemma}

\begin{proof}\leanok
    \uses{thm:sector_bnt_bijective_match_of_sameMPV,
        lem:sector_bnt_coeff_identity_via_matched_mpv_phase,
        lem:sector_bnt_copy_weight_matching_of_coeff_identity}
    The full-basis matching theorem gives the bijection and block gauges. The
    self-overlap comparison
    \begin{align}
        \langle V_N(Q_k),V_N(Q_k)\rangle
        &=
        |\zeta_k|^{2N}
        \langle V_N(P_{\beta(k)}),V_N(P_{\beta(k)})\rangle
        \notag
    \end{align}
    together with the BNT normalization limits on both sides forces
    $|\zeta_k|=1$. The matched MPV phase identities also give the eventual
    coefficient identities
    $c_N^{(\beta(k))}(P)=\zeta_k^N c_N^{(k)}(Q)$.
    Lemma~\ref{lem:sector_bnt_copy_weight_matching_of_coeff_identity} applies
    the finite power-sum comparison in each sector and gives the copy-weight
    identities.
\end{proof}

\begin{theorem}[Full-basis global gauge in matched coordinates]%
    \label{thm:sector_bnt_equal_global_gauge}
    \lean{MPSTensor.ft_sector_bnt_equal_global_gauge}
    \leanok
    \uses{def:sector_bnt_cf}
    Let $P$ and $Q$ be BNT canonical forms whose total tensors have the same
    MPV family at every positive length. Then there exist:
    \begin{itemize}
        \item a bijection $\beta:J(Q)\simeq J(P)$,
        \item bond-dimension equalities
              $\dim P_{\beta(k)}=\dim Q_k$,
        \item copy-number equalities $r_{\beta(k)}(P)=r_k(Q)$ and
              copy permutations $\tau_k$,
        \item unit-modulus phases $\zeta_k$ with $|\zeta_k|=1$ and block
              gauges $X_k$,
    \end{itemize}
    such that, for every sector $k$, copy $q$, and physical index $i$,
    \begin{align}
        Q_k^i
        &= \zeta_k\,X_k P_{\beta(k)}^i X_k^{-1},
        \notag\\
        \nu_{k,q}
        &= \zeta_k^{-1} \mu_{\beta(k),\tau_k(q)}.
        \notag
    \end{align}
    Moreover, if $X=\bigoplus_k (\mathbf{1}_{r_k}\otimes X_k)$ is the flattened
    block-diagonal gauge, then, for every physical index~$i$, the flattened
    $Q$-coordinates satisfy
    \begin{align}
        Q_{\mathrm{flat}}^i
        &= X\,P_{\mathrm{matched}}^i\,X^{-1}.
        \notag
    \end{align}
    This is the explicit direct-sum gauge construction of
    \cite[Appendix MPV proof, lines~1189--1192]{Cirac2016MPDO_arXiv}.
    % Per-sector unit-modulus restriction eliminated by the exact matcher;
    % closure recorded in
    % docs/paper-gaps/cpsv16_global_vs_persector_unit_witness.tex.
\end{theorem}

\begin{proof}\leanok
    \uses{lem:sector_bnt_equal_matched_copy_weight_witnesses,
        thm:sector_bnt_global_gauge_of_copy_weight_matching}
    Lemma~\ref{lem:sector_bnt_equal_matched_copy_weight_witnesses} gives
    $\beta$, the bond-dimension identifications, the phases $\zeta_k$, the
    block gauges $X_k$, and copy permutations $\tau_k$ satisfying
    $\nu_{k,q}=\zeta_k^{-1}\mu_{\beta(k),\tau_k(q)}$.
    Theorem~\ref{thm:sector_bnt_global_gauge_of_copy_weight_matching} combines
    the block conjugacies into the flattened direct-sum conjugation by
    $X=\bigoplus_k(\mathbf{1}_{r_k}\otimes X_k)$.
\end{proof}

\begin{theorem}[Literal equal-MPV gauge after sector-coordinate permutation]%
    \label{thm:sector_bnt_equal_mps_gaugeEquiv_literal}
    \lean{MPSTensor.ft_sector_bnt_equal_mps_gaugeEquiv_literal}
    \leanok
    \uses{def:sector_bnt_cf}
    Let $P$ and $Q$ be BNT canonical forms whose total tensors have the same
    MPV family at every positive length. Then the total bond dimensions agree;
    write their common value as $D$. There is
    \begin{align}
        Y
        &\in \GL_D(\C)
        \label{eq:ft_literal_gauge_Y}
    \end{align}
    such that, for every physical index $i$,
    \begin{align}
        Q_{\mathrm{tot}}^i
        &= Y\,P_{\mathrm{tot}}^i\,Y^{-1},
        \notag
    \end{align}
    where $P_{\mathrm{tot}}^i$ and $Q_{\mathrm{tot}}^i$ are both regarded as
    matrices in $\MN{D}$ using the total bond-dimension equality. This is the
    BNT form of
    \cite[Corollary~II.2 and Appendix MPV proof, lines~1189--1192]
    {Cirac2016MPDO_arXiv}.
\end{theorem}

\begin{proof}\leanok
    \uses{thm:sector_bnt_equal_global_gauge}
    Apply Theorem~\ref{thm:sector_bnt_equal_global_gauge} to obtain the
    flattened gauge after the sectors of $P$ have been reindexed by the
    bijection between basis blocks of $P$ and $Q$. The sector-coordinate
    permutation $\rho$, induced by this bijection, the copy permutations, and
    the bond-dimension equalities, identifies the reindexed $P$-side direct sum
    with the total tensor $P_{\mathrm{tot}}$, transported along the total
    bond-dimension equality. Composing the permutation gauge for $\rho$ with
    the sector-permuted gauge gives the single matrix $Y$ of
    (\ref{eq:ft_literal_gauge_Y}).
\end{proof}

\begin{lemma}[Unitary gauge in matched sector coordinates]
    \label{lem:sector_bnt_equal_unitary_matched_gauge}
    \lean{MPSTensor.ft_sector_bnt_equal_unitary_global_gauge_witnesses}
    \leanok
    \uses{thm:sector_bnt_equal_global_gauge,
        lem:left_canonical_irreducible_same_phase_unitary_gauge,
        lem:unitary_global_gauge_of_blocks}
    Let $P$ and $Q$ be BNT canonical forms whose total tensors generate the
    same MPV family at every positive length. There are a sector bijection
    $\beta$, copy bijections $\tau_k$, phases $\zeta_k$, sector unitaries $U_k$,
    and a unitary block-diagonal matrix $X$ such that
    \begin{align}
        Q_k^i
        &= \zeta_k U_kP_{\beta(k)}^iU_k^\dagger,
        \notag\\
        \nu_{k,q}
        &= \zeta_k^{-1}\mu_{\beta(k),\tau_k(q)},
        \notag\\
        |\zeta_k|
        &= 1.
        \notag
    \end{align}
    The matched sectors and copies have equal bond dimensions and
    multiplicities. If $P_{\mathrm{match}}$ denotes the corresponding
    reordered block assembly, then
    $Q_{\mathrm{flat}}^i=XP_{\mathrm{match}}^iX^{-1}$.
\end{lemma}
\begin{proof}\leanok
    \uses{lem:sector_bnt_equal_matched_copy_weight_witnesses,
        lem:left_canonical_irreducible_same_phase_unitary_gauge,
        lem:unitary_global_gauge_of_blocks}
    Apply Lemma~\ref{lem:left_canonical_irreducible_same_phase_unitary_gauge}
    to each matched sector, retaining the phase in the copy-weight identity.
    Repeating $U_k$ on every copy gives $X=\bigoplus_{k,q}U_k$, which is unitary
    by Lemma~\ref{lem:unitary_global_gauge_of_blocks}. For every matched copy
    $(k,q)$,
    \begin{align}
        \nu_{k,q}Q_k^i
        &= \zeta_k^{-1}\mu_{\beta(k),\tau_k(q)}
          \zeta_k U_kP_{\beta(k)}^iU_k^\dagger
        = \mu_{\beta(k),\tau_k(q)}U_kP_{\beta(k)}^iU_k^\dagger.
        \notag
    \end{align}
    These are precisely the blocks of
    $Q_{\mathrm{flat}}^i=XP_{\mathrm{match}}^iX^{-1}$.
\end{proof}

\begin{corollary}[Canonical Form II unitary refinement in the equal case]%
    \label{cor:sector_bnt_equal_unitary_global_gauge}
    \lean{MPSTensor.fundamentalTheorem_equal_canonicalForm_unitary}
    \leanok
    \uses{lem:sector_bnt_equal_unitary_matched_gauge}
    Let $P$ and $Q$ be BNT canonical forms whose total tensors generate the
    same MPV family at every positive length. Their total bond dimensions
    agree; writing the common value as $D$, there is a unitary matrix
    $U\in\mathrm{U}(D)$ such that, for every physical index~$i$,
    \begin{align}
        Q_{\mathrm{tot}}^i
        &= U P_{\mathrm{tot}}^i U^\dagger.
        \notag
    \end{align}
    This is the equal case of
    \cite[Corollary~C.5]{Cirac2016MPDO_arXiv}.
\end{corollary}

\begin{proof}\leanok
    \uses{lem:sector_bnt_equal_unitary_matched_gauge}
    For each matched sector, retain the phase in both identities
    \begin{align}
        Q_k^i
        &= \zeta_k U_kP_{\beta(k)}^iU_k^\dagger,
        \notag\\
        \nu_{k,q}
        &= \zeta_k^{-1}\mu_{\beta(k),\tau_k(q)}.
        \notag
    \end{align}
    For every matched copy $(k,q)$,
    \begin{align}
        \nu_{k,q}\zeta_k
        &= \zeta_k^{-1}\mu_{\beta(k),\tau_k(q)}\zeta_k
        = \mu_{\beta(k),\tau_k(q)}.
        \notag
    \end{align}
    Hence $Q_{\mathrm{flat}}^i=XP_{\mathrm{match}}^iX^{-1}$. Repeating $U_k$
    on its matched copies gives
    \begin{align}
        X
        &= \bigoplus_{k,q}U_k,
        \notag\\
        X^\dagger X
        &= \mathbb 1.
        \notag
    \end{align}
    If $\Pi$ is the permutation matrix restoring the original sector and copy
    order, then $\Pi^\dagger\Pi=\mathbb 1$. Hence
    \begin{align}
        U
        &= X\Pi,
        \notag\\
        U^\dagger U
        &= \mathbb 1,
        \notag
    \end{align}
    and the permutation of matched coordinates gives
    \begin{align}
        Q_{\mathrm{tot}}^i
        &= UP_{\mathrm{tot}}^iU^\dagger.
        \notag
    \end{align}
\end{proof}
